% GENERATED FILE. Edit canonical lessons, then run npm run generate:books.
% canonical-source-hash: fnv1a64:141e2950045b5615
% canonical-lessons: FA-C235-tiz, FA-C235-mojud, FA-C235-kafi, FA-C235-majani, FA-C235-bikar

\chapter{Sharp, Available, Enough, Free of charge, Unemployed}
\label{ch:fa-sharp-available-enough-free-of-charge-unemployed}
\hlchaptermodality{235}

\begin{chapteropening}
\textbf{By the end of this chapter:} \emph{I can say sharp, available, enough, free of charge and unemployed.}

Complete the last lesson of chapter 235: I can say sharp, available, enough, free of charge and unemployed.
\end{chapteropening}

\section[\texorpdfstring{\fa{تیز} (tiz)}{تیز (tiz)}]{\fa{تیز} (tiz) — sharp}

\label{lesson:FA-C235-tiz}



\begin{quote}
Before the new one: say the Persian for life, then the Persian for a part.
\end{quote}

\subsection*{You'll want to know: \fa{تیز}}
\textbf{\fa{تیز}} — \emph{tiz} — \textquotedblleft{}sharp\textquotedblright{}.

\begin{cousinweb}
\emph{tiz} is a native Persian word, inherited rather than borrowed; it is usually traced to Middle Persian \emph{tēz}, \textquotedblleft{}sharp\textquotedblright{}. EtymDB takes it back to the Indo-European root \emph{*(s)teyg-}, \textquotedblleft{}to be sharp, to prick\textquotedblright{}, the root usually also given for English \emph{stick}.
\end{cousinweb}

\begin{grammarlens}[title={What you've built}]
One more word for this chapter.
\end{grammarlens}

\subsection*{Your turn}
\begin{itemize}
\raggedright
  \item \emph{Say it:} \emph{tiz}
  \item \emph{Say it:} \emph{tiz}, once more
  \item \emph{Say it:} say \emph{tiz}
  \item \emph{Recall:} say the Persian for life, then the Persian for a part, then say \emph{tiz} again
\end{itemize}

\begin{tcolorbox}[breakable,colback=teal!4,colframe=teal!35!black,title={Before you move on}]
What does \fa{تیز} mean? (\textquotedblleft{}Sharp\textquotedblright{}.) Say it once more.
\end{tcolorbox}
\section[\texorpdfstring{\fa{موجود} (mojud)}{موجود (mojud)}]{\fa{موجود} (mojud) — available, in stock}

\label{lesson:FA-C235-mojud}



\begin{quote}
Before the new one: say the Persian for a part, then the Persian for sharp.
\end{quote}

\subsection*{You'll want to know: \fa{موجود}}
\textbf{\fa{موجود}} — \emph{mojud} — \textquotedblleft{}available, in stock\textquotedblright{}.

\begin{cousinweb}
\emph{mojud} is borrowed from Arabic \emph{mawjūd}, \textquotedblleft{}found, existing\textquotedblright{}. Its Arabic root is \emph{w-j-d} (Lane).
\end{cousinweb}

\begin{grammarlens}[title={What you've built}]
One more word for this chapter.
\end{grammarlens}

\subsection*{Your turn}
\begin{itemize}
\raggedright
  \item \emph{Say it:} \emph{mojud}
  \item \emph{Say it:} \emph{mojud}, once more
  \item \emph{Say it:} say \emph{mojud}
  \item \emph{Recall:} say the Persian for a part, then the Persian for sharp, then say \emph{mojud} again
\end{itemize}

\begin{tcolorbox}[breakable,colback=teal!4,colframe=teal!35!black,title={Before you move on}]
What does \fa{موجود} mean? (\textquotedblleft{}Available, in stock\textquotedblright{}.) Say it once more.
\end{tcolorbox}
\section[\texorpdfstring{\fa{کافی} (kâfi)}{کافی (kâfi)}]{\fa{کافی} (kâfi) — enough}

\label{lesson:FA-C235-kafi}



\begin{quote}
Before the new one: say the Persian for sharp, then the Persian for available.
\end{quote}

\subsection*{You'll want to know: \fa{کافی}}
\textbf{\fa{کافی}} — \emph{kâfi} — \textquotedblleft{}enough\textquotedblright{}.

\begin{cousinweb}
\emph{kâfi} is borrowed from Arabic \emph{kāfī}, \textquotedblleft{}sufficient\textquotedblright{}. Its Arabic root is \emph{k-f-y} (Lane).
\end{cousinweb}

\begin{grammarlens}[title={What you've built}]
One more word for this chapter.
\end{grammarlens}

\subsection*{Your turn}
\begin{itemize}
\raggedright
  \item \emph{Say it:} \emph{kâfi}
  \item \emph{Say it:} \emph{kâfi}, once more
  \item \emph{Say it:} say \emph{kâfi}
  \item \emph{Recall:} say the Persian for sharp, then the Persian for available, then say \emph{kâfi} again
\end{itemize}

\begin{tcolorbox}[breakable,colback=teal!4,colframe=teal!35!black,title={Before you move on}]
What does \fa{کافی} mean? (\textquotedblleft{}Enough\textquotedblright{}.) Say it once more.
\end{tcolorbox}
\section[\texorpdfstring{\fa{مجانی} (majâni)}{مجانی (majâni)}]{\fa{مجانی} (majâni) — free of charge}

\label{lesson:FA-C235-majani}



\begin{quote}
Before the new one: say the Persian for available, then the Persian for enough.
\end{quote}

\subsection*{You'll want to know: \fa{مجانی}}
\textbf{\fa{مجانی}} — \emph{majâni} — \textquotedblleft{}free of charge\textquotedblright{}.

\begin{cousinweb}
\emph{majâni} is borrowed from Arabic \emph{majjānī}, \textquotedblleft{}free, without payment\textquotedblright{}. Its Arabic root is \emph{m-j-n} (Lane).
\end{cousinweb}

\begin{grammarlens}[title={What you've built}]
One more word for this chapter.
\end{grammarlens}

\subsection*{Your turn}
\begin{itemize}
\raggedright
  \item \emph{Say it:} \emph{majâni}
  \item \emph{Say it:} \emph{majâni}, once more
  \item \emph{Say it:} say \emph{majâni}
  \item \emph{Recall:} say the Persian for available, then the Persian for enough, then say \emph{majâni} again
\end{itemize}

\begin{tcolorbox}[breakable,colback=teal!4,colframe=teal!35!black,title={Before you move on}]
What does \fa{مجانی} mean? (\textquotedblleft{}Free of charge\textquotedblright{}.) Say it once more.
\end{tcolorbox}
\section[\texorpdfstring{\fa{بیکار} (bikâr)}{بیکار (bikâr)}]{\fa{بیکار} (bikâr) — unemployed}

\label{lesson:FA-C235-bikar}



\begin{quote}
Before the new one: say the Persian for enough, then the Persian for free of charge.
\end{quote}

\subsection*{You'll want to know: \fa{بیکار}}
\textbf{\fa{بیکار}} — \emph{bikâr} — \textquotedblleft{}unemployed\textquotedblright{}.

\begin{cousinweb}
\emph{bikâr} is \emph{bi-}, \textquotedblleft{}without\textquotedblright{}, with \emph{kâr}, \textquotedblleft{}work\textquotedblright{}: without work. Both are native Persian; \emph{kâr} continues Middle Persian \emph{kār} (EtymDB). (EtymDB also lists Arabic \emph{bikr}, \textquotedblleft{}virgin\textquotedblright{}, a different word.)
\end{cousinweb}

\begin{grammarlens}[title={What you've built}]
One more word for this chapter.
\end{grammarlens}

\subsection*{Your turn}
\begin{itemize}
\raggedright
  \item \emph{Say it:} \emph{bikâr}
  \item \emph{Say it:} \emph{bikâr}, once more
  \item \emph{Say it:} say \emph{bikâr}
  \item \emph{Recall:} say the Persian for enough, then the Persian for free of charge, then say \emph{bikâr} again
\end{itemize}

\begin{tcolorbox}[breakable,colback=teal!4,colframe=teal!35!black,title={Before you move on}]
What does \fa{بیکار} mean? (\textquotedblleft{}Unemployed\textquotedblright{}.) Say it once more.
\end{tcolorbox}
